% !TEX root = paper.tex
% =============================================================
% RAR -- Conclusion and Limitations v0
% Source-of-truth: project_context.md (binding constraints).
% Stage 3 of paper-writing skill. Draft 0 intro frozen 2026-07-09.
% Last updated: 2026-07-09
%
% Three disclosed limitations are LOCKED 2026-07-09 :
%   1. few-shot regimes
%   2. reconstruction-branch latency
%   3. high old/new class similarity (worst cases = ISCX-Tor + EBSNN-WebApp)
%
% Plus two explicit non-claims:
%   - adversarial evasion of the reconstruction branch
%   - extreme concept drift outside the encrypted-traffic domain
%
% Conventions:
%   - KDD venue: integrated related work, colon-style subtitles,
%     no \smartparagraph, no em-dashes.
%   - Headings are claims. Topic sentences assert claims.
%   - Disclosure, not omission: limitations stated openly.
% =============================================================

\section{Conclusion}\label{sec:conclusion}

RAR transforms per-sample reconstruction errors into knowledge ownership probabilities, enabling adaptive soft routing between old and new classification heads. By introducing reconstruction-aware routing, dual-head decoupling, and logit calibration, RAR effectively mitigates knowledge interference in encrypted traffic class-incremental learning. Extensive experiments on six encrypted-traffic datasets demonstrate that RAR consistently reduces forgetting while improving mixed-class recognition performance over existing CIL baselines.

Although effective, RAR still relies on limited replay samples to characterize old knowledge and may face challenges when old and new classes exhibit highly similar representations. Future work will explore more adaptive knowledge routing mechanisms and integrate reconstruction-aware routing with incremental representation learning to further enhance continual traffic classification.

